\section*{Chapitre \ref{ch:pl:post-trade-systems} --- Systèmes de post-marché}

\begin{solution}{exo:pl:post-trade-systems:1}
Les clés déterminent quelle confirmation est comparée à quelle transaction~; une tolérance à ce niveau associerait une transaction à la confirmation d'une autre et produirait alors
un faux écart ou une fausse concordance. Les champs comparés sont ceux où un écart porte une information~--- et seul le prix a une raison
légitime (la précision) de différer légèrement.
\end{solution}

\begin{solution}{exo:pl:post-trade-systems:2}
Un arrondi à quatre décimales déplace le prix d'au plus une demi-unité à la quatrième décimale, soit 0.00005~\$~; à cette tolérance, toute confirmation correcte
concorde, ce qui marque le début du palier de la figure du chapitre.
\end{solution}

\begin{solution}{exo:pl:post-trade-systems:3}
91 des 138 écarts ont été corrigés pendant les fenêtres du soir~; 47 ont échoué, soit 34~\%.
\end{solution}

\begin{solution}{exo:pl:post-trade-systems:4}
Un mauvais compte sur une confirmation provoque un écart sur une seule transaction, que le \omterm{def:pl:post-trade-systems:recon}{rapprochement} détecte. Un mauvais compte dans la table des SSI est recopié dans chaque
instruction destinée à cette contrepartie et à ce marché, concorde des deux côtés parce que les deux l'utilisent, et fait échouer toutes les transactions concernées jusqu'à ce que l'agent
les rejette.
\end{solution}

\begin{solution}{exo:pl:post-trade-systems:5}
Par un regroupement un-à-plusieurs~: la clé du dépositaire est associée aux deux affectations, dont les quantités sont additionnées pour donner 3\,000 avant la comparaison, de sorte que
les deux côtés concordent et qu'aucun écart n'est levé.
\end{solution}

\begin{solution}{exo:pl:post-trade-systems:6}
Parce que seules les transactions dont le règlement est prévu pour vendredi figurent dans les positions de vendredi. Sous T+2, les transactions de jeudi et de vendredi se règlent la semaine suivante, et 6 des
10 échecs les concernaient~; les 4 autres ont chacun provoqué un écart de position.
\end{solution}

\begin{solution}{exo:pl:post-trade-systems:7}
47, 35, 22, 14, 11 et 10 échecs pour 3 à 8 personnes~: il en faut huit pour atteindre exactement les 10 échecs de T+2, sept pour s'en approcher à un près.
\end{solution}

\begin{solution}{exo:pl:post-trade-systems:8}
Le point de pourcentage supplémentaire représente cinq prix erronés désormais acceptés comme concordances~: les écarts n'ont pas disparu, ils sont restés masqués jusqu'à ce que
les montants en espèces ne concordent plus. Un taux de concordance n'a de sens qu'à une tolérance dérivée de la précision des contreparties, et doit être présenté avec le nombre
d'écarts relevés ensuite par le \omterm{def:pl:post-trade-systems:recon}{rapprochement}.
\end{solution}

\begin{solution}{pb:pl:post-trade-systems:1}
\begin{enumerate}
\item La sécurité, le sens, la contrepartie, la date de transaction et le fonds sont des clés et doivent être exacts~; le prix, la quantité, la date de règlement et le compte sont comparés,
le prix à l'intérieur d'une tolérance.
\item L'accord du gestionnaire d'investissement sur la justesse des conditions, ce qui permet à la transaction de se régler sans intervention supplémentaire.
\item Ils arrondissent le prix moyen à quatre décimales, de sorte que leurs prix corrects diffèrent de ceux de la société d'au plus 0.00005~\$~; à 0.0000001~\$, 60.4~\%
des transactions concordent.
\item De 0.00005~\$ à 0.01~\$, à 95.8~\%.
\item Cinq prix erronés sont acceptés, et le taux monte à 96.8~\%.
\item Une livraison contre paiement de 9\,400 actions pour 917\,317.25~\$ à la date de règlement~: la sécurité, la quantité, le montant et la date proviennent de la transaction,
le compte du fonds provient du modèle canonique, l'agent et le compte du courtier de la table des SSI.
\item Proposée, approuvée par une seconde personne, valable à compter d'une date, jamais modifiée sur place.
\item Les clés~: les regroupements un-à-plusieurs et les positions agrégées, ainsi que les identifiants qui diffèrent entre les deux côtés.
\item Un type, un propriétaire et un âge.
\item Un écart qui survit à plusieurs \omterm{def:pl:post-trade-systems:recon}{rapprochements} est un écart que personne ne comprend, et c'est là que s'accumulent les pertes et les erreurs d'information.
\item Le jour ouvré suivant (08~h~00--17~h~00) sous T+2, le soir de la date de transaction (17~h~00--21~h~00) sous T+1~: les heures limites citées dans le chapitre.
\item Environ 26 heures par jour~; la journée donne trois personnes 27, la soirée 12.
\item 10 transactions, 2.74~millions de dollars, sous T+2~; 47, 19.17~millions de dollars, sous T+1.
\item Les mauvais comptes et les confirmations manquantes, 18 chacun.
\item 4 écarts de position, tous vieux de un à trois jours, sous T+2~; 24 sous T+1, dont 7 de moins d'un jour, 12 de un à trois jours et 5 de trois à sept jours.
\item \textbf{Résultat nommé.} À une tolérance de prix de 0.0000001~\$, 60.4~\% des transactions concordent automatiquement~; de 0.00005~\$ à 0.01~\$, 95.8~\%, à
0.05~\$, 96.8~\% avec cinq prix erronés acceptés. Sur les 138 écarts de la semaine, 10 restent non corrigés à l'heure limite sous un cycle de deux jours, et 47 sous un
cycle d'un jour, ce qui fait échouer des transactions pour 2.74~millions et 19.17~millions de dollars.
\item Sept personnes ramènent les échecs à 11, huit à 10.
\item Les échecs tombent à 13 avec les mêmes trois personnes.
\item Diviser les erreurs par deux~: cela coûte un travail sur les données de référence et les flux une seule fois, au lieu de quatre salaires supplémentaires chaque année, et réduit chaque
coût ultérieur des écarts, pas seulement celui des échecs.
\item Cela raccourcit le temps nécessaire pour corriger un écart, sans changer leur nombre.
\end{enumerate}
\end{solution}

\begin{solution}{iq:pl:post-trade-systems:1}
Un élément sur lequel deux enregistrements qui devraient concorder diffèrent. Le rapport donne la clé, les deux valeurs, le type, le propriétaire, la date d'ouverture et sa
tranche d'âge, classé par valeur et par âge.
\iqlookfor{Type, propriétaire, âge~; valeur d'abord.}
\end{solution}

\begin{solution}{iq:pl:post-trade-systems:2}
Clés exactes~; tolérance nulle sur les quantités, les dates et les comptes~; tolérance de prix par contrepartie, fondée sur la précision qu'elle utilise~; puis vérifier le taux de
concordance par rapport aux écarts relevés ensuite, pour s'assurer que la tolérance ne masque pas d'erreurs.
\iqlookfor{Tolérance fondée sur la précision, par contrepartie.}
\end{solution}

\begin{solution}{iq:pl:post-trade-systems:3}
Associer chaque position de fonds et de stratégie au compte du dépositaire grâce aux données de référence, additionner les positions de la société dans les clés du dépositaire
(un-à-plusieurs), comparer, et lever des écarts assortis de propriétaires~; conserver l'association versionnée pour que le \omterm{def:pl:post-trade-systems:recon}{rapprochement} de la veille puisse être relancé.
\iqlookfor{Association des clés et regroupements un-à-plusieurs.}
\end{solution}

\begin{solution}{iq:pl:post-trade-systems:4}
L'affectation, la confirmation et l'\omterm{def:pl:post-trade-systems:confirmation}{affirmation} passent à la date de transaction~; le temps de correction d'un écart se réduit à quelques heures, de sorte que moins d'écarts (SSI, affectation
anticipée) comptent davantage que des personnes supplémentaires~; le \omterm{def:pl:position-and-pnl-services:service}{service de positions} doit suivre les positions réglées et non réglées, et le financement avance d'un jour.
\iqlookfor{Moins d'écarts, plus tôt.}
\end{solution}

\begin{solution}{iq:pl:post-trade-systems:5}
La transaction canonique issue de l'enregistrement~; un moteur d'appariement (matching engine) à tolérances par contrepartie~; un suivi des \omterm{def:pl:post-trade-systems:confirmation}{affirmations} par rapport à l'heure limite~; un référentiel de SSI approuvées,
datées, qui génèrent les instructions~; des \omterm{def:pl:post-trade-systems:recon}{rapprochements} des positions, des espèces et des transactions, avec appariement un-à-plusieurs et agrégé~; une file d'écarts unique
avec types, propriétaires et vieillissement~; un taux de \omterm{def:pl:trade-capture-and-booking:stp}{traitement de bout en bout} par contrepartie comme mesure.
\iqlookfor{Une file d'écarts unique, mesurée en amont.}
\end{solution}
